\chapter{LITERATURE REVIEW}\label{ch:lit}

This chapter reviews the work on which the project builds and against which it is compared. It covers seven themes:
the clinical need for seizure detection, wearable and multimodal detectors, the SeizeIT programme and its datasets,
detection algorithms, evaluation practice, networks built from lookup tables, and hardware implementations of seizure
detectors. The survey covers 79 references. Each was matched against its published record, and numbers are quoted as
reported in the cited abstracts. Each theme ends with a summary table, and Section~\ref{sec:lit-gaps} lists the
research gaps that the project addresses.

\section{Epilepsy and the Clinical Need for Seizure Detection}\label{sec:lit-clinical}

Epilepsy is one of the most common chronic neurological disorders: the Global Burden of Disease study estimated 45.9
million people with active epilepsy in 2016, an age-standardised prevalence of 621.5 per 100,000
\cite{gbd2019epilepsy}. The International League Against Epilepsy (ILAE) classifies seizures by their onset (focal,
generalised or unknown) and, for focal seizures, by whether awareness is kept \cite{fisher2017ilae}. This distinction
matters for detection: a focal seizure with impaired awareness may produce little movement, while a focal-to-bilateral
tonic-clonic seizure produces strong muscle activity.

Seizure counts are the basis of treatment decisions, yet they are unreliable. Elger and Hoppe identify lack of
awareness that a seizure has occurred as the primary cause of under-reporting, and call for ambulatory long-term
detection that combines heart-rate and movement measures \cite{elger2018underreporting}. Unattended seizures are also
dangerous: in the MORTEMUS study of 147 epilepsy monitoring units, the incidence of sudden unexpected death in epilepsy
(SUDEP) was 5.1 per 1000 patient-years, and in the recorded cases terminal cardiorespiratory dysfunction followed a
generalised tonic-clonic seizure \cite{ryvlin2013mortemus}. The ILAE--IFCN clinical practice guideline therefore
recommends clinically validated wearable devices for tonic-clonic seizures, especially in unsupervised patients, but
does not recommend their clinical use for other seizure types, for which the evidence is insufficient
\cite{beniczky2021guideline}. Detecting focal seizures with a wearable, the target of this project, is therefore an
open clinical problem.

\section{Wearable and Multimodal Seizure Detection}\label{sec:lit-wearable}

Most validated wearables detect convulsive seizures from movement or muscle activity. A wearable surface-EMG device on
the arm detected tonic-clonic seizures with a sensitivity above 90\,\% and a latency within 30\,s in a prospective
three-centre study of 71 patients \cite{beniczky2018emg}. Empatica wristbands combining electrodermal activity and
accelerometry reached a sensitivity of 94.55\,\% at 0.2 false alarms per day for convulsive seizures in 69 patients
\cite{onorati2017empatica}. A review of multimodal detection found sensitivities between 4 and 100\,\% and 0.25 to 20
false detections per 8\,h, with the best results for tonic-clonic seizures \cite{leijten2018review}, and a 2022
meta-analysis of 28 studies pooled a tonic-clonic sensitivity of 0.91 at 2.1 false alarms per 24\,h while reporting
focal-seizure sensitivities from 31.1 to 93.1\,\% \cite{naganur2022meta}. These devices work for convulsive seizures,
but movement and skin signals carry little information about non-motor focal seizures.

Brain signals are needed for focal seizures, and behind-the-ear (bte) EEG makes them wearable. Vandecasteele et al.
showed that two bte channels allow visual seizure annotation (sensitivity 65.7\,\%) and automated detection, with a
patient-independent detector reaching 64.1\,\% at 2.8 false positives per 24\,h \cite{vandecasteele2020bte}. Adding
ECG to bte-EEG raised the sensitivity by 11 and 8 percentage points at the same false-alarm rate in two of three
databases \cite{vandecasteele2021ecg}, and in a semi-automated setting ECG cut the false detections from 6.5 to 2.4
per day \cite{bhagubai2023ecg}. Fusion has the same effect for tonic-clonic seizures: on 27 SeizeIT2 participants, EEG
alone reached 100\,\% sensitivity at 10.3 false positives per 24\,h, while EEG, EMG and accelerometry together reached
90.9\,\% at 0.1 \cite{zhang2024tcs}. Missing sensors are rarely handled; the only seizure work found on this problem is
a cyclic transformer that learns from paired EEG and EMG and is tested on EEG alone \cite{zhang2024cyclic}. The
closest edge-oriented multimodal detector, BrainFuseNet, fuses four EEG channels with photoplethysmography and
accelerometry on a microcontroller and reaches 64.28\,\% sample sensitivity at 0.21 false positives per hour
\cite{ingolfsson2024brainfusenet}; it uses multiply-accumulate arithmetic and no ECG or EMG. Table~\ref{tab:lit-wearable}
summarises these studies.

\begin{table}[H]
\centering\small
\caption{Wearable and multimodal seizure detection studies (results as reported)}\label{tab:lit-wearable}
\begin{tabularx}{\textwidth}{|P{30mm}|c|P{22mm}|P{22mm}|L|L|}
\toprule
\textbf{Work} & \textbf{Year} & \textbf{Signals} & \textbf{Seizures} & \textbf{Result} & \textbf{Gap for our problem}\\
\midrule
Beniczky et al. \cite{beniczky2018emg} & 2018 & EMG & tonic-clonic & $>$90\,\% sensitivity & convulsive only\\ \hline
Onorati et al. \cite{onorati2017empatica} & 2017 & EDA + ACC & convulsive & 94.55\,\%, 0.2 FA/day & convulsive only\\ \hline
Naganur et al. \cite{naganur2022meta} & 2022 & meta-analysis & tonic-clonic & 0.91, 2.1 FA/24\,h & focal not solved\\ \hline
Vandecasteele et al. \cite{vandecasteele2020bte} & 2020 & bte-EEG & focal & 64.1\,\%, 2.8 FP/24\,h (patient-independent) & EEG only\\ \hline
Vandecasteele et al. \cite{vandecasteele2021ecg} & 2021 & bte-EEG + ECG & focal & +8--11\,\% sensitivity with ECG & software only\\ \hline
Bhagubai et al. \cite{bhagubai2024challenge} & 2024 & bte-EEG & challenge & tree ensemble won & no hardware budget\\ \hline
Zhang et al. \cite{zhang2024tcs} & 2024 & EEG + EMG + ACC & tonic-clonic & 90.9\,\%, 0.1 FP/24\,h & tonic-clonic only\\ \hline
Swinnen et al. \cite{swinnen2026seizeit2} & 2026 & SeizeIT2 wearable & focal, 192 adults & 0.73 at precision 0.004 & false alarms\\ \hline
BrainFuseNet \cite{ingolfsson2024brainfusenet} & 2024 & EEG + PPG + ACC & -- & 64.28\,\%, 0.21 FP/h, MCU & multipliers, no ECG/EMG\\ \hline
Zhang et al. \cite{zhang2024cyclic} & 2024 & EEG + EMG & -- & -- & large transformer\\ \hline

\end{tabularx}
\end{table}

\section{The SeizeIT Programme and the SeizeIT2 Dataset}\label{sec:lit-seizeit}

The SeizeIT projects record patients with the Sensor Dot wearable (Byteflies), which carries bte-EEG, ECG, EMG, an
accelerometer and a gyroscope. SeizeIT1 (42 patients with focal epilepsy) was the training set of the 2023 ICASSP
wearable-EEG grand challenge, which was evaluated on recordings from the wearable itself; the winning method was a
feature-based tree ensemble, and the organisers named the lack of training data as the main obstacle
\cite{bhagubai2024challenge}. The winning entry, SeizFt, used hand-crafted features, Fourier-transform-surrogate
augmentation and a CatBoost ensemble \cite{alhussaini2023seizft}.

SeizeIT2 \cite{bhagubai2025seizeit2} is the first open dataset of wearable recordings from patients with focal
epilepsy: 125 patients from five European epilepsy monitoring units, more than 11,000 hours and 883 focal seizures
annotated on video-EEG, with a proposed training and validation split, benchmark methods (an SVM and ChronoNet) and
evaluation metrics. It is the dataset of this project. A multicentre validation of the same device in 192 adults with
616 focal seizures shows why false alarms are the main problem: the algorithm alone reached a sensitivity of 0.73 but
a precision of only 0.004, and expert review was needed to raise the precision to 0.83 \cite{swinnen2026seizeit2}.
Other SeizeIT2 studies use ECG-only anomaly detection \cite{reintjes2025ecg} or a complex-valued multimodal network
\cite{su2026cmepinet}; none of them reports a hardware implementation.

\section{Seizure Detection Algorithms}\label{sec:lit-alg}

\textbf{Feature-based models.} Classical detectors compute features such as band power and line length and feed them
to a classifier. Shoeb and Guttag's patient-specific SVM detected 96\,\% of 173 test seizures with a median of two false
detections per 24\,h on the CHB-MIT scalp dataset, but it needs seizures of the same patient for training
\cite{shoeb2010icml}. On SeizeIT data, tree ensembles have been the strongest: besides SeizFt
\cite{alhussaini2023seizft}, a random forest on band-power features reached a sensitivity of 73.6\,\% when trained on
SeizeIT1 and tested on SeizeIT2 \cite{huang2023lightweight}.

\textbf{Deep models.} ChronoNet stacks multi-scale one-dimensional convolutions and gated recurrent units and was
designed for abnormal-EEG identification \cite{roy2019chrononet}; it is one of the two SeizeIT2 baselines. Trained on
SeizeIT1, it reached only 15.2\,\% sensitivity on SeizeIT2 data \cite{huang2023lightweight}. Compact convolutional
networks such as EEGNet \cite{lawhern2018eegnet} and EEGWaveNet \cite{thuwajit2022eegwavenet}, graph networks
\cite{tang2022gnnssl} and transformers \cite{wu2025seizuretransformer,yang2023biot} improve results on scalp EEG, but
they rely on full electrode montages, large models and floating-point multiplications, which do not fit a
battery-powered two-channel wearable.

\textbf{Imperfect labels.} Seizure annotations have uncertain onsets and offsets. Weak supervision on clinical
annotations \cite{saab2020weak}, attention-based multiple-instance learning \cite{isaev2020mil}, noise-aware training
\cite{shama2026bundl} and distillation from a large teacher into a few-channel student \cite{zheng2023kdgnn} address
this, but only for floating-point models. They motivate our distillation and noise-robust training experiments
(extensions N4 and N5). Table~\ref{tab:lit-alg} summarises the algorithms.

\begin{table}[H]
\centering\small
\caption{Seizure detection algorithms and evaluation frameworks (results as reported)}\label{tab:lit-alg}
\begin{tabularx}{\textwidth}{|P{30mm}|c|P{25mm}|P{20mm}|L|L|}
\toprule
\textbf{Work} & \textbf{Year} & \textbf{Method} & \textbf{Data} & \textbf{Result} & \textbf{Gap for our problem}\\
\midrule
Shoeb \& Guttag \cite{shoeb2010icml} & 2010 & patient-specific SVM & CHB-MIT & 96\,\% of seizures, 2 FD/24\,h & patient-specific, scalp\\ \hline
SeizFt \cite{alhussaini2023seizft} & 2023 & features + CatBoost & SeizeIT1/2 & ICASSP 2023 winner & EEG only, no hardware cost\\ \hline
Huang et al. \cite{huang2023lightweight} & 2023 & random forest; ChronoNet & SeizeIT1/2 & RF 73.6\,\%; ChronoNet 15.2\,\% & EEG only\\ \hline
ChronoNet \cite{roy2019chrononet} & 2019 & CNN + GRU & TUH abnormal & benchmark at the time & large, multipliers\\ \hline
EEGWaveNet \cite{thuwajit2022eegwavenet} & 2022 & multiscale CNN & CHB-MIT, TUSZ & better subject-independent & scalp, float CNN\\ \hline
Tang et al. \cite{tang2022gnnssl} & 2022 & self-supervised GNN & TUSZ & AUROC 0.875 & needs full montage\\ \hline
Shama \& Venkataraman \cite{shama2026bundl} & 2026 & noisy-label training & TUH, CHB-MIT & robust to label noise & scalp EEG\\ \hline
Zheng et al. \cite{zheng2023kdgnn} & 2023 & distilled GNN & TUSZ & competitive with 2 channels & float student\\ \hline
SzCORE \cite{dan2024szcore} & 2024 & evaluation framework & scalp EEG & standard event metrics & not for wearables\\ \hline
SzCORE challenge \cite{dan2025szcorechallenge} & 2025 & blind benchmark & 65 subjects & best event-F1 43\,\% & scalp EEG\\ \hline

\end{tabularx}
\end{table}

\section{Evaluation Practice}\label{sec:lit-eval}

Published numbers are hard to compare because scoring differs: epoch-based and any-overlap scoring give different
results for the same system \cite{ziyabari2017metrics}, and many works report sample-level sensitivity or specificity
instead of detected seizures and false alarms per day. The SzCORE framework standardises event-based scoring for
seizure detection \cite{dan2024szcore}. In its first blind benchmark on 65 subjects and 4,360 hours of EEG, the best
event-F1 was 43\,\%, and the organisers reported a gap between published and blind-test performance
\cite{dan2025szcorechallenge}. A second source of optimism is data leakage: most surveyed deep-learning EEG studies mix
data of the same subject between training and test sets, which can greatly overestimate performance on new subjects
\cite{brookshire2024leakage}. Large public corpora such as TUSZ \cite{shah2018tusz} are scalp recordings, not
wearable ones. This project therefore uses SzCORE event scoring, patient-wise cross-validation and a test set used only
once.

\section{Lookup-Table Native, Weightless and Logic Networks}\label{sec:lit-lut}

\textbf{Weightless networks} replace weighted sums with memory lookups. WiSARD stored binary patterns in RAM nodes
without gradient training \cite{aleksander1984wisard}; BTHOWeN and ULEEN made weightless networks smaller and more
energy-efficient for edge inference \cite{susskind2022bthowen,susskind2023uleen}. The Differentiable Weightless Neural
Network (DWN) made both the table contents and the input mapping trainable by gradient descent and reported better
latency, throughput, energy and area than earlier models on FPGAs and microcontrollers \cite{bacellar2024dwn}. DWNs
have since been applied to inertial activity recognition at 56--104\,nJ per sample \cite{bacellar2025nanoml} and to
continuous control \cite{kresse2025dwc}. The DWN is the model family used in this project.

\textbf{Logic-gate networks} learn a network of two-input Boolean gates; the discretised networks classify more than a
million images per second on one CPU core \cite{petersen2022dlgn}, with convolutional, lighter and recurrent variants
\cite{petersen2024convdlgn,ruttgers2026lightdlgn,buhrer2025rdlgn}. \textbf{LUT networks for FPGAs}, including LUTNet,
LogicNets, PolyLUT, NeuraLUT, AmigoLUT and ReducedLUT, map sparse quantised neurons onto FPGA lookup tables and reach
latencies of tens of nanoseconds \cite{wang2019lutnet,umuroglu2020logicnets,andronic2023polylut,andronic2024neuralut,weng2025amigolut,cassidy2025reducedlut};
TreeLUT maps boosted decision trees onto LUTs with a 4--101$\times$ lower area-delay product than recent LUT networks
\cite{khataei2025treelut}. Mecik and Kumm showed that the thermometer encoder of a DWN can raise its LUT count by up to
3.2$\times$ and dominates the cost of small networks \cite{mecik2025thermo}, a cost most papers leave out.

These networks have been tested mostly on images, tabular data and particle physics. The closest biosignal works are
inter-patient ECG arrhythmia classification with logic-gate and LUT networks \cite{mommen2026ecg} and logic-gate
networks for EEG classification on an embedded CPU \cite{dharia2026eeglgn}. The only weightless seizure work found is a
WiSARD study of seizure anticipation in four patients, reported as accuracy only \cite{deaguiar2015wisard}.
Table~\ref{tab:lit-lut} summarises this family.

\begin{table}[H]
\centering\small
\caption{LUT-native, weightless and logic networks (results as reported)}\label{tab:lit-lut}
\begin{tabularx}{\textwidth}{|P{30mm}|c|P{27mm}|P{22mm}|L|L|}
\toprule
\textbf{Work} & \textbf{Year} & \textbf{Model} & \textbf{Domain} & \textbf{Hardware / result} & \textbf{Gap}\\
\midrule
WiSARD \cite{aleksander1984wisard} & 1984 & RAM-node WNN & images & -- & no gradient training\\ \hline
ULEEN \cite{susskind2023uleen} & 2023 & weightless & MLPerf Tiny & 7.1$\times$ energy eff.\ vs FINN & no biosignals\\ \hline
DWN \cite{bacellar2024dwn} & 2024 & differentiable LUT6 net & tabular & FPGA, MCU & stateless, no events\\ \hline
nanoML HAR \cite{bacellar2025nanoml} & 2025 & DWN & inertial HAR & 56--104\,nJ / sample & single modality\\ \hline
LGN \cite{petersen2022dlgn} & 2022 & 2-input logic gates & images & $>$1\,M images/s (CPU) & no biosignals\\ \hline
Recurrent LGN \cite{buhrer2025rdlgn} & 2025 & stateful logic gates & translation & -- & no hardware\\ \hline
LogicNets \cite{umuroglu2020logicnets} & 2020 & sparse quantised LUT net & jet tagging & $<$15\,ns latency & tiny fan-in\\ \hline
NeuraLUT \cite{andronic2024neuralut} & 2024 & dense sub-nets in LUTs & jet tagging & 4.3$\times$ lower latency & static\\ \hline
TreeLUT \cite{khataei2025treelut} & 2025 & quantised GBDT & tabular & 4--101$\times$ lower area-delay & per window\\ \hline
Mecik \& Kumm \cite{mecik2025thermo} & 2025 & DWN generator & jet tagging & encoder up to 3.2$\times$ LUTs & --\\ \hline
Mommen et al. \cite{mommen2026ecg} & 2026 & LGN, LUT network & ECG arrhythmia & 2.0--3.0k LUTs, 5--7\,mW & single modality\\ \hline
Dharia et al. \cite{dharia2026eeglgn} & 2026 & LGN & EEG & embedded CPU & no FPGA, no seizures\\ \hline
de Aguiar et al. \cite{deaguiar2015wisard} & 2015 & WiSARD & seizure anticipation & -- & 4 patients, accuracy only\\ \hline

\end{tabularx}
\end{table}

\section{Hardware Implementations of Seizure Detectors}\label{sec:lit-hw}

\textbf{ASICs} give the lowest energy per decision. Verma et al. integrated EEG acquisition with spectral feature
extraction at 9\,\textmu J per feature vector \cite{verma2010micro}; patient-specific SVM classifiers reached
2.73 and 1.83\,\textmu J per classification \cite{altaf2015jssc,altaf2016tbcas}; boosted decision trees for
intracranial EEG reached 41.2\,nJ per classification \cite{shoaran2018jetcas}; and NeuralTree classifies 256 channels
at 0.227\,\textmu J per classification \cite{shin2022neuraltree}. \textbf{FPGA} detectors include an STFT and SVM
design on a Zynq at 380\,mW \cite{wang2018access}, on-chip SVM training \cite{elhosary2019tbcas}, a low-power
patient-specific design for intracranial EEG at 0.49\,mW dynamic power \cite{razi2022tbcas}, further FPGA designs
\cite{vittimberga2025sensors}, and, closest to our work, a spiking neural network on a small FPGA that needs 4.55\,nJ per
inference and reaches 0.3 false positives per hour on CHB-MIT \cite{busia2025fpga}. \textbf{Microcontroller} systems
include e-Glass, which runs for 2.71 days on one charge \cite{sopic2018eglass}, the EEGformer transformer at 0.31\,mJ
per inference \cite{busia2024eegformer}, BrainFuseNet \cite{ingolfsson2024brainfusenet} and low-power EEG platforms
\cite{kartsch2019biowolf,ingolfsson2022artifact}. Generic tools for efficient inference, such as binarised networks,
FINN, integer-only quantisation and hls4ml \cite{hubara2016bnn,umuroglu2017finn,jacob2018quant,duarte2018hls4ml},
reduce but do not remove multiplications.

Across these works, almost all detectors use EEG alone, many are trained per patient, and results are reported as
epoch or sample accuracy rather than detected seizures and false alarms per day. None of the surveyed detectors builds
its classifier from FPGA lookup tables, and none of them reports how the classifier hardware is checked for faults
during use. Table~\ref{tab:lit-hw} summarises them.

\begin{table}[H]
\centering\footnotesize
\caption{Hardware seizure detectors (results as reported; units differ between papers)}\label{tab:lit-hw}
\begin{tabularx}{\textwidth}{|P{25mm}|@{\hspace{4pt}}c@{\hspace{4pt}}|P{17mm}|P{16mm}|P{22mm}|L|L|}
\toprule
\textbf{Work} & \textbf{Year} & \textbf{Platform} & \textbf{Signals} & \textbf{Hardware} & \textbf{Detection} & \textbf{Gap}\\
\midrule
Verma et al. \cite{verma2010micro} & 2010 & ASIC & scalp EEG & 9\,\textmu J / feature vector & classifier off-chip & EEG only\\ \hline
Bin Altaf et al. \cite{altaf2015jssc} & 2015 & ASIC & scalp EEG & 2.73\,\textmu J / class. & 95.7\,\% (epoch) & per patient\\ \hline
Shoaran et al. \cite{shoaran2018jetcas} & 2018 & ASIC & iEEG & 41.2\,nJ / class. & F1 87.9\,\% & implant\\ \hline
NeuralTree \cite{shin2022neuraltree} & 2022 & ASIC & EEG, iEEG & 0.227\,\textmu J / class. & 95.6\,\% (EEG) & neural only\\ \hline
Wang et al. \cite{wang2018access} & 2018 & FPGA (Zynq) & scalp EEG & 380\,mW & 98.4\,\%, 0.356 FD/h & per patient\\ \hline
Farhang Razi \& Schmid \cite{razi2022tbcas} & 2022 & FPGA (Cyclone V) & iEEG & 0.49\,mW dynamic & 100\,\% sens. & implant\\ \hline
Busia et al. \cite{busia2025fpga} & 2025 & FPGA (iCE40) & scalp EEG & 4.55\,nJ / inference & 0.3 FP/h & EEG only\\ \hline
e-Glass \cite{sopic2018eglass} & 2018 & MCU & scalp EEG & 2.71 days / charge & 93.8\,\% (epoch) & EEG only\\ \hline
EEGformer \cite{busia2024eegformer} & 2024 & MCU & scalp EEG & 0.31\,mJ / inference & 73\,\%, 0.15 FP/h & multipliers\\ \hline
BrainFuseNet \cite{ingolfsson2024brainfusenet} & 2024 & MCU & EEG, PPG, ACC & 0.11\,mJ / inference & 0.21 FP/h & no ECG/EMG\\ \hline

\end{tabularx}
\end{table}

\section{Research Gaps}\label{sec:lit-gaps}

The survey shows that (i) focal seizures remain the open clinical problem, and false alarms are what stops wearable
detectors from being used; (ii) multimodal fusion is the most consistent way to cut false alarms, but hardware
detectors use EEG alone; (iii) LUT-native networks are among the most efficient models for FPGAs but have never been
used for seizure detection; and (iv) evaluation practice varies so much that results are rarely comparable.
Table~\ref{tab:lit-gaps} lists the specific gaps and how this project addresses them.

\begin{xltabular}{\textwidth}{|c|P{45mm}|L|L|}
\caption{Research gaps and how this project addresses them}\label{tab:lit-gaps}\\
\toprule
\textbf{\#} & \textbf{Gap} & \textbf{Evidence} & \textbf{How this project addresses it}\\
\midrule
\endfirsthead
\toprule
\textbf{\#} & \textbf{Gap} & \textbf{Evidence} & \textbf{How this project addresses it}\\
\midrule
\endhead
\small 1 & \small No LUT-native (weightless, logic-gate or LUT) seizure detector exists & \small LUT networks are benchmarked on images, physics and tabular data; the only weightless seizure work is a 4-patient study \cite{deaguiar2015wisard} & \small First benchmark of DWN, LGN, TreeLUT and LogicNets-style networks for seizure detection, on SeizeIT2\\ \hline
\small 2 & \small Hardware detectors use EEG alone, although fusion is the most consistent way to cut false alarms & \small \cite{vandecasteele2021ecg,zhang2024tcs}; Section~\ref{sec:lit-hw} & \small Integer features from behind-the-ear EEG, ECG, EMG and motion in one bit-exact front end\\ \hline
\small 3 & \small Missing sensors are rarely handled; LUT networks never & \small \cite{zhang2024cyclic}; Section~\ref{sec:lit-lut} & \small Presence bits, modality dropout, gated and late fusion, missing-sensor tests (N6)\\ \hline
\small 4 & \small The encoder's hardware cost is ignored, although it dominates small LUT networks & \small \cite{mecik2025thermo}; our estimate: about 850 of 960 LUTs & \small Cost model that counts every comparator; comparator-budgeted and learned encoders (N2)\\ \hline
\small 5 & \small LUT networks are stateless; seizures unfold over tens of seconds & \small \cite{buhrer2025rdlgn} & \small A stateful DWN with state bits fed back every second (N1)\\ \hline
\small 6 & \small Results are patient-specific, leak between training and test, or use incomparable metrics & \small \cite{brookshire2024leakage,ziyabari2017metrics,dan2025szcorechallenge} & \small Patient-wise cross-validation, a locked test set used once, SzCORE metrics with the paper's metrics alongside\\ \hline
\small 7 & \small Small models learn from noisy labels and from larger models only in floating point & \small \cite{shama2026bundl,zheng2023kdgnn} & \small Noise-robust and multiple-instance losses (N5) and distillation from LightGBM (N4) into a LUT network\\ \hline
\small 8 & \small Detectors are black boxes & \small Section~\ref{sec:lit-lut} & \small Tracing alarms to the inputs read by each LUT (N7)\\ \hline
\small 9 & \small Published wearable baselines are not independently reproduced & \small Section~\ref{sec:lit-wearable} & \small Independent reproduction of the SeizeIT2 ChronoNet baseline, with and without test-patient information\\ \hline

\end{xltabular}
